\chapter*{附录A}
\begin{mybox}{方向余弦矩阵 (DCM) 的七大性质}
对于任何旋转矩阵 $R \in SO(3)$，其数学性质如下：

\begin{enumerate}
    \item \textbf{逆矩阵等于转置 (Orthogonality)}: 
    根据定义 $R^T R = I$，可以直接导出：
    \begin{equation}
        R^{-1} = R^T
    \end{equation}
    这意味着求旋转矩阵的逆运算只需要进行矩阵转置。

    \item \textbf{行列式约束 (Unimodular)}: 
    $\det(R) = 1$。这保证了变换是刚性的（不改变体积且不含镜像反射）。若 $\det(R) = -1$，则表示包含镜像的坐标系，不属于右手系。

    \item \textbf{保持向量长度 (Isometry)}: 
    对任意向量 $x \in \mathbb{R}^3$，旋转后的向量 $y = Rx$ 模长不变：
    \begin{equation}
        \|Rx\| = \sqrt{(Rx)^T (Rx)} = \sqrt{x^T R^T R x} = \sqrt{x^T I x} = \|x\|
    \end{equation}

    \item \textbf{保持内积与夹角}: 
    旋转不改变两个向量之间的点积：
    \begin{equation}
        (Rx) \cdot (Ry) = (Rx)^T (Ry) = x^T R^T R y = x^T y = x \cdot y
    \end{equation}
    因此，刚体上各点间的相对夹角在旋转前后保持不变。

    \item \textbf{保持外积 (Cross Product Conservation)}: 
    旋转满足对叉乘的分配律：
    \begin{equation}
        R(x \times y) = (Rx) \times (Ry)
    \end{equation}
    这说明旋转矩阵不仅作用于向量，也保持了向量间的几何拓扑关系。

    \item \textbf{矩阵的列向量与行向量的几何意义}: 
    \begin{itemize}
        \item $R$ 的三列是物体坐标系 $\{b\}$ 的轴在固定系 $\{s\}$ 中的坐标表达。
        \item $R$ 的三行是固定坐标系 $\{s\}$ 的轴在物体系 $\{b\}$ 中的坐标表达。
    \end{itemize}

    \item \textbf{群乘法封闭性 (Closure)}: 
    若 $R_1, R_2 \in SO(3)$，则它们的乘积 $R_3 = R_1 R_2$ 依然满足 $R_3^T R_3 = I$ 且 $\det(R_3) = 1$，即 $R_3 \in SO(3)$。这构成了连续旋转的复合基础。
\end{enumerate}
\end{mybox}